%% =====================================================================
%%  Corrigés — Unit 1 « Hello! »
%% =====================================================================
\corrigesUnite{1}{Hello!}

\begin{corrige}{1}
\begin{reponses}[2]
\item 8 h 15 : \rep{Good morning!}
\item 14 h 30 : \rep{Good afternoon!}
\item 19 h 00 : \rep{Good evening!}
\item 23 h 00 : \rep{Good night!}
\end{reponses}
\pourquoi{\en{good morning} s'emploie jusqu'à midi, \en{good afternoon} de midi à 18 heures
environ, \en{good evening} ensuite. À 23 heures, Léa \textbf{part} se coucher : c'est donc
\en{good night}, qui se dit uniquement pour prendre congé.}
\piege{dire \en{Good night!} en arrivant à une soirée. Pour « bonsoir » en arrivant, c'est
toujours \en{Good evening!} — ou tout simplement \en{Hello!}, qui convient à toute heure.}
\end{corrige}

\begin{corrige}{2}
\begin{reponses}[5]
\item \rep{d}
\item \rep{e}
\item \rep{b}
\item \rep{a}
\item \rep{c}
\end{reponses}
\pourquoi{à \en{Nice to meet you} on répond en reprenant la formule avec \en{too} (aussi) ; à un
remerciement, \en{You're welcome} (« de rien ») ; à \en{See you tomorrow!}, une formule
d'au revoir.}
\retiens{\en{You're welcome} ne veut pas dire « bienvenue » mais « de rien ». « Bienvenue ! »
se dit \en{Welcome!}, sans \en{you're}.}
\end{corrige}

\begin{corrige}{3}
\begin{reponses}[5]
\item \rep{he}
\item \rep{she}
\item \rep{he}
\item \rep{she}
\item \rep{it}
\item \rep{it}
\item \rep{they}
\item \rep{we}
\item \rep{you}
\item \rep{they}
\end{reponses}
\pourquoi{\en{the studio} et \en{New York} ne sont pas des personnes : on emploie \en{it}.
\en{Léa and I} = « Léa et moi » = « nous » (\en{we}). \en{you and Jack} = « toi et Jack » =
« vous » (\en{you}). Deux femmes (\en{Emma and Mrs.\ Lopez}) : \en{they}, comme pour deux hommes.}
\piege{utiliser \en{he} ou \en{she} pour une ville ou un objet, comme en français
(« le studio ? Il est… »). En anglais, une chose est toujours \en{it}.}
\end{corrige}

\begin{corrige}{4}
\begin{reponses}[4]
\item \rep{am}
\item \rep{are}
\item \rep{is}
\item \rep{is}
\item \rep{is}
\item \rep{are}
\item \rep{are}
\item \rep{is}
\end{reponses}
\pourquoi{\en{my name} et \en{it} sont au singulier : \en{is}. \en{Jack and Emma} sont deux
personnes (= \en{they}) : \en{are}.}
\retiens{\en{I am} · \en{he / she / it is} · \en{you / we / they are}.}
\end{corrige}

\begin{corrige}{5}
\begin{reponses}[2]
\item \rep{I'm Léa.}
\item \rep{You're welcome.}
\item \rep{He's fine.}
\item \rep{She's Emma.}
\item \rep{We're OK.}
\item \rep{They're in New York.}
\item \rep{What's your name?}
\item \rep{It's nice to meet you.}
\end{reponses}
\pourquoi{l'apostrophe remplace la première lettre du verbe : \en{a} de \en{am} et de \en{are},
\en{i} de \en{is}.}
\piege{écrire \en{Your welcome}. \en{your} veut dire « ton, ta, votre » ; ici, il faut
\en{you're} = \en{you are}.}
\end{corrige}

\begin{corrige}{6}
\begin{reponses}[2]
\item \rep{My name is Jack.}
\item \rep{Nice to meet you.}
\item \rep{How are you?}
\item \rep{I'm fine, thanks.}
\item \rep{What's your name?}
\item \rep{See you tomorrow!}
\end{reponses}
\pourquoi{en anglais comme en français, la phrase commence par une majuscule et le point
d'interrogation se place à la fin. En revanche, on ne met \textbf{pas d'espace} avant
\en{?} et \en{!} en anglais.}
\end{corrige}

\begin{corrige}{7}
\begin{reponses}[3]
\item \rep{My} name is Léa.
\item \rep{I'm} Jack.
\item What's \rep{your} name?
\item \rep{You're} welcome.
\item \rep{I'm} fine, thanks.
\item \rep{You're} very nice!
\end{reponses}
\pourquoi{devant un nom (\en{name}), on met un possessif : \en{my} (mon), \en{your} (ton).
Devant un prénom ou un adjectif (\en{fine}, \en{nice}), il faut un sujet et le verbe \en{be} :
\en{I'm}, \en{you're}.}
\retiens{\en{I'm} = « je suis » ; \en{my} = « mon, ma ». Astuce : remplacez par la forme
pleine. Si \en{I am} ou \en{you are} fonctionne, c'est la forme avec apostrophe.}
\end{corrige}

\begin{corrige}{8}
\begin{reponses}[4]
\item \rep{hi}
\item \rep{hand}
\item \rep{air}
\item \rep{heat}
\item \rep{old}
\item \rep{heart}
\item \rep{ear}
\item \rep{his}
\end{reponses}
\pourquoi{quand vous entendez un léger souffle avant la voyelle, c'est le mot avec \en{h}.
Si vous avez fait plus de deux erreurs, réécoutez la piste~1.09 en regardant la bouche
dans un miroir, puis refaites l'exercice. Le texte de la piste est page~\pageref{trans-1.10}.}
\end{corrige}

\begin{corrige}{9}
\begin{center}\otTableau
\begin{tabular}{@{}l@{\hspace{8mm}}l@{\hspace{8mm}}l@{\hspace{8mm}}l@{}}
 & \otTH{Registre} & \otTH{Bonjour ou au revoir ?} & \otTH{Les indices}\\\midrule
{\otSemi A} & informel & bonjour   & \en{Hey}, \en{Hi}, des prénoms\\
{\otSemi B} & formel   & bonjour   & \en{Good afternoon}, \en{Mrs.\ Lopez}, \en{Nice to meet you}\\
{\otSemi C} & informel & au revoir & \en{Bye}, \en{See you later}\\
{\otSemi D} & formel   & au revoir & \en{Goodbye}, \en{Mrs.\ Lopez}, \en{Have a nice day}\\\bottomrule
\end{tabular}\end{center}
\retiens{les indices du registre formel : \en{good morning / afternoon / evening}, \en{goodbye},
\en{Mr.} ou \en{Mrs.} + nom de famille. Les indices du registre informel : \en{hi}, \en{hey},
\en{bye}, les prénoms.}
\end{corrige}

\begin{corrige}{10}
\begin{anglais}\otTableau
\begin{tabular}{@{}l@{\hspace{6mm}}l@{}}
\otNomDit{sam} & Hi! \rep{I'm} Sam.\\
\otNomDit{lea} & Hello, Sam. My \rep{name} is Léa.\\
\otNomDit{sam} & \rep{Nice} to meet you, Léa.\\
\otNomDit{lea} & Nice to meet you \rep{too}.\\
\otNomDit{sam} & How \rep{are} you?\\
\otNomDit{lea} & I'm \rep{fine}, thanks. And you?\\
\otNomDit{sam} & I'm OK. Have a nice \rep{day}!\\
\otNomDit{lea} & Thank you. You \rep{too}!\\
\end{tabular}
\end{anglais}
\piege{confondre \en{too} (aussi) et \en{to} (« à », dans \en{Nice to meet you}) : ces deux mots
se prononcent presque pareil, mais \en{too} est accentué, \phon{tuː}.}
\end{corrige}

\begin{corrige}{11}
\begin{reponses}[1]
\item \rep{Good morning. My name is Paul.} \textit{ou} \rep{Good morning. I'm Paul.}
\item \rep{Nice to meet you!}
\item \rep{How are you? — I'm very well, thank you. And you?}
\item \rep{Bye, Emma! See you tomorrow!}
\item \rep{Good evening, Mrs.\ Lopez.}
\item \rep{Thank you very much. — You're welcome.}
\item \rep{This is Jack.}
\item \rep{They're in New York.} \textit{(ou} \rep{They are in New York.}\textit{)}
\end{reponses}
\pourquoi{phrase 1, le matin : \en{good morning} plutôt que \en{hello}, qui reste possible.
Phrase 4 : « salut » se dit pour arriver (\en{hi}) \textbf{et} pour partir (\en{bye}) ; ici Léa
part, donc \en{bye}. Phrase 5 : en arrivant le soir, \en{good evening} et non \en{good night}.}
\piege{traduire mot à mot « je m'appelle » par \en{I call me}. On dit \en{My name is…} ou
\en{I'm…}.}
\end{corrige}

\begin{corrige}{12}
Réponses modèles (avec votre prénom à la place de Léa) :
\begin{reponses}[1]
\item \rep{Hello, Jack! I'm Léa.} \textit{ou} \rep{Hi! My name is Léa.}
\item \rep{Nice to meet you too.}
\item \rep{I'm fine, thanks. And you?} \textit{ou} \rep{I'm good, thanks. And you?}
\item \rep{Thank you!} \textit{ou} \rep{Thanks!}
\item \rep{Bye, Jack! Have a nice evening!} \textit{ou} \rep{Goodbye! Have a nice evening!}
\end{reponses}
\retiens{refaites l'exercice deux ou trois fois, jusqu'à répondre sans hésiter pendant la
pause. Parler à voix haute, même seul, est la meilleure façon de fixer les phrases.}
\end{corrige}

\begin{corrigetest}
\begin{reponses}[5][french]
\item \rep{a}
\item \rep{c}
\item \rep{b}
\item \rep{b}
\item \rep{a}
\item \rep{c}
\item \rep{b}
\item \rep{c}
\item \rep{a}
\item \rep{c}
\end{reponses}
\pourquoi{question 4 : \en{She is Emma}, contracté en \en{She's Emma}. Question 8 : dans
\en{hour}, le \en{h} est muet, comme dans le français « heure ». Question 10 : \en{Good night}
se dit seulement pour prendre congé.}
\end{corrigetest}
